\scanpage{134}
% archive foliation: 133
donc à toutes les structures possibles du type envisagé sur $X$) sont en correspondance biunivoque, via restriction, avec l'ens. des prolongements $p_3$ de $p_2$, à l'ordre 3 (sans autre condition sur ceux-ci).
Ce sont \inserted{donc} les sections sur $X$ d'un torseur canonique sur $S$, dans le faisceau abélien
\[\underline{\operatorname{Hom}}(\underline\omega,\underline\omega^{\otimes3})\simeq\underline\omega^{\otimes2}\]
\[
\deleted{\unsure{p\in}}\quad\underbrace{\mathcal C_{X/S}}_{\substack{\text{ensemble des}\\\text{``connexions}\\\text{projectives''}}}\quad\text{torseur sous }\underline\omega^{\otimes2}=\Omega^{1\,\otimes2}_{X/S}.
\]
On dit aussi que les structures du type envisagé sur $X$ sont des \uline{connexions projectives} sur $X$, et le torseur envisagé est le \uline{torseur des connexions projectives} sur $X/S$. [Elles sont aussi en correspondance biunivoque avec l'ens. des connexions \inserted{correspondantes} $c$ sur $P=P(E)$ ($E$ comme dessus), ou ce qui revient au même, des connexions \inserted{(respectant les L.)} sur le fibré en Alg. de Lie associé (image \deleted{par}\inserted{directe} $P\longrightarrow X$ de $\mathcal T_{P/X}$, ou encore $E\otimes\check E/\mathcal O_X$) --- on trouve
\[\unsure{\mathcal H}\simeq(\operatorname{Sym}^2E)\otimes\underline\omega^{-1}\]
extension successive --- en commençant par le bas) de $\underline\omega^1$, $\underline\omega^0$, $\underline\omega^{-1}$ --- et il faut prendre les connexions $\nabla$
\[
c:\unsure{\mathcal H}\longrightarrow\unsure{\mathcal H}\otimes\underline\omega\simeq\operatorname{Sym}^2E
\]
qui, sur $\underline\omega\subset\unsure{\mathcal H}$ (applique nécessairement dans $\operatorname{Ker}(\operatorname{Sym}^2E\longrightarrow\operatorname{Sym}^2\mathcal O_X\simeq\mathcal O_X)$), ext. de $\underline\omega$ par $\underline\omega^2$,
